\cleardoublepage
\chapter{TRANSKRIP WAWANCARA PENGUJIAN KEGUNAAN}
\label{chap:lampiran-b}

Bagian ini memuat transkrip verbatim sesi \textit{task-based usability test} dan wawancara \textit{debriefing} terhadap kelima responden pengujian kegunaan sebagaimana dibahas pada Subbab \ref{chap:evaluasi}. Kode responden (R1-R5) dan profil masing-masing mengikuti Tabel \ref{tab:profil-responden}. Beberapa kata pengisi (\textit{filler words}) dan pengulangan lisan dipertahankan apa adanya untuk menjaga keaslian data.

\section{Transkrip Wawancara Responden 1 (Muhammad Zufar Dafy)}

\subsubsection*{1. Pendahuluan \& Persiapan}

\textbf{Farah:} Halo, terima kasih sudah bersedia meluangkan waktu Anda. Nama saya Farah Aulia, Anda bisa memanggil saya Farah. Saat ini saya sedang menguji sebuah aplikasi, yaitu teman belajar algoritma dan struktur data sebagai bagian dari tugas akhir saya. Nah, yang perlu digarisbawahi di sini yang kita akan uji adalah aplikasinya, bukan kamu. Jadi kalau nanti ada bagian yang membingungkan atau sulit ditemukan, itu adalah masukan bagi saya, jadi bukan kesalahan kamu. Kemudian sesi ini akan berlangsung kurang lebih sekitar 30 menit, terdiri dari kamu mengerjakan beberapa tugas di aplikasi, lalu mengisi kuesioner singkat. Selama mengerjakan tugas, saya meminta kamu untuk berpikir keras atau \textit{think aloud}. Jadi coba ucapkan apa yang kamu pikirkan atau cari saat menggunakan aplikasi. Misalnya, "Saya cari tombol ini di sini karena biasanya ada di pojok kanan atas," atau "Saya bingung nih, harus klik yang mana." Enggak perlu dijawab dengan benar, saya cuma perlu tahu apa yang kamu pikirkan. Saya akan mengamati dan mencatat, tapi saya tidak akan membantu selama pengerjaan tugas. Itu memang dari bagian dari pengujian. Kalau kamu benar-benar \textit{stuck}, kamu boleh bilang "Saya menyerah" dan kita lanjut ke tugas berikutnya. Sebelum itu, apakah ada pertanyaan untuk memulai?

\textbf{Zufar:} Oke, cukup ya, Mbak.

\textbf{Farah:} Oke, terima kasih. Baik karena saya rasa tidak ada pertanyaan, kita bisa masuk ke sesi pengujian. Saya izin untuk mengambil HP saya terlebih dahulu. Oke. Selanjutnya coba kamu \textit{share screen} ke \textit{link} yang sudah saya berikan sebelumnya terkait \textit{website} saya.

\bigskip

\subsubsection*{2. Sesi Pengujian Aplikasi (Usability Testing)}

\textbf{Zufar:} Kayak gini, Mbak?

\textbf{Farah:} Bentar ya, Mas. Sepertinya nge-\textit{lag}. Oke, ya benar seperti itu, Mas. Nah, untuk tugas pertama, silakan daftar terlebih dahulu.

\textbf{Zufar:} Ini biasanya buat tombol daftar tuh pasti di bawah tombol \textit{login} ya masuk umumnya. Di sini jelas, terus... nama lengkap...

\textbf{Farah:} Oke, baik. Terima kasih. Jadi saya izin menjelaskan, ini adalah halaman utama ketika Anda selesai melakukan pendaftaran. Berikutnya, coba lakukan untuk mencari topik \textit{linked list} di \textit{website} ini. Saya mulai dari sekarang.

\textbf{Zufar:} Cari topik ya... Oke, ini buat navigasinya biasanya di atas ya emang paling gampang dilihat juga. Oke. Udah, udah dapet topiknya.

\textbf{Farah:} Oh iya, silakan masuk ke \textit{linked list} ya.

\textbf{Zufar:} Gimana, Mbak?

\textbf{Farah:} \textit{Linked list}, \textit{linked list}.

\textbf{Zufar:} Oh, \textit{list} linear itu.

\textbf{Farah:} Oke, baik. Setelah itu, coba baca halaman materi yang ditampilkan lalu tandai sebagai selesai.

\textbf{Zufar:} Ini saya langsung ceritain baca aja dulu deh. Oke buat tombol selesainya di sini ya, kelihatan kok Mbak.

\textbf{Farah:} Oke, terima kasih. Berikutnya masuk ke contoh dan baca bagian contohnya. Setelah itu tekan tombol selesai apabila sudah. Selanjutnya kita masuk ke latihan. Di sini latihan ada soal dari AI dan soal sendiri. Jadi soal dari AI ini, AI akan men-\textit{generate} soal dan akan menampilkan kepada Anda, dan dia akan mengevaluasi secara otomatis jawaban Anda. Nah, itu mungkin bisa dikerjakan dan berikan... dan kemudian baca evaluasi yang diberikan. Bisa kerjakan soal satu aja terlebih dahulu.

\textbf{Zufar:} Ini kalau lihat ke belakang lagi bisa nggak ya, Mas? \textit{(Melihat layar)}

\textbf{Farah:} Jawab ngawur aja nggak apa-apa, aku tidak menilai kamu bisa atau tidak. Sebisa kamu aja.

\textbf{Zufar:} Oke, ini misal saya jawab... pasti kurang paham gitu nah.

\textbf{Farah:} Hahaha, gila kamu ngasih kurang paham. Bisa-bisa ya ngasih kurang paham. Mending kamu ngawur dulu kayak misal nih, "Jelaskan perbedaan antara \textit{list} linear biasa..." kalau \textit{list} linear biasa mungkin kayak pointernya cuman satu, kalau \textit{double linked list} dia bisa ke belakang ke depan gitu.

\textbf{Zufar:} Biasa pointernya satu, sedangkan \textit{double list} pointernya dua. Untuk...

\textbf{Farah:} Aduh aku harus ngedit video ini anjir gara-gara omongan kita hahaha.

\textbf{Zufar:} \textit{Double linked list} lebih diuntungkan ketika... diperlukan perbandingan antar \textit{value} dalam \textit{list} dan perlu komputasi cepat. Terus gimana nih? Nilai soal ini ya Mbak?

\textbf{Farah:} Iya. Iya benar.

\textbf{Zufar:} Jawaban saya betul atau tidak. Wah nilainya 55. "Penjelasan lebih mendalam tentang perbedaan kedua struktur..." Oh...

\textbf{Farah:} Oke, setelah itu coba tekan tandai selesai dan masuk ke... Setelah itu tekan tandai selesai dan coba Anda buka \textit{tab} ringkasan dan selesaikan alur belajar untuk topik \textit{list} linear ini.

\textbf{Zufar:} Oke, berarti ini beresin dulu atau langsung ringkasan aja?

\textbf{Farah:} Ringkasan aja. Dan tandai selesai. Oke, baik, terima kasih. Oke terakhir, coba akses halaman progres belajar dan cek status penyelesaian topik \textit{linked list}.

\textbf{Zufar:} Oke ini buatnya tadi kan kita pindah... pindah menu lagi ya berarti kita ke navigasi atas, progres. Ini ya, jelas. Terus tadi apa ya Mbak?

\textbf{Farah:} \textit{Link list}. \textit{List} linear.

\textbf{Zufar:} Linear ya. Tadi udah ya ditandain.

\textbf{Farah:} Ya, kita sudah melihat bahwa \textit{list} linear telah selesai dipelajari. Oke saya rasa itu adalah semua tugas sudah selesai.

\bigskip

\subsubsection*{3. Sesi Kuesioner (SUS) \& Debriefing}

\textbf{Farah:} Sekarang saya akan kirimkan \textit{link} Google Form berisi 10 pernyataan singkat tentang pengalamanmu menggunakan aplikasi ini. Isi berdasarkan kesan pertamamu, nggak perlu dipikir terlalu lama, dan mohon diisi dulu sebelum kita ngobrol supaya jawabannya murni dari pengalamanmu sendiri, bukan hasil diskusi kita. Itu tadi \textit{link}-nya sudah saya kirim di WhatsApp yang Google Form tadi, silakan diisi.

\textbf{Zufar:} Ini udah saya buka ya \textit{link}-nya. Ini saya perlu ngeluarin pikiran saya juga ya saat ngisi ya?

\textbf{Farah:} Ah tidak perlu, tidak perlu. Silakan isi sesuai pengalaman Anda.

\textit{(Jeda saat Zufar mengisi kuesioner)}

\textbf{Farah:} Oke, baik, terima kasih setelah mengisi kuesioner SUS-nya. Oke, selanjutnya kita masuk ke \textit{debriefing}. Secara umum bagaimana kesan kamu memakai aplikasi ini?

\textbf{Zufar:} Buat aplikasi ini cukup jelas ya buat tujuan aplikasinya apa gitu. Jadi secara UI itu berarti udah menggambarkan banget lah ya si aplikasi ini tentang apa. Terus \textit{user friendly} juga, ini udah mengikuti standar juga menurut saya, maksudnya udah seperti aplikasi pada umumnya ya jadi pembiasaannya nggak terlalu banyak gitu. Kita udah langsung familiar gitu ya. Seperti itu. Terus mungkin cukup kaget ya sama fiturnya, soalnya bagus gitu Mbak.

\textbf{Farah:} Oh makasih, makasih Mas. Oke lanjut, apa menurutmu bagian yang paling membingungkan atau susah digunakan? Dan apa sarannya?

\textbf{Zufar:} Susah digunakan ya? Menurut saya sih udah cukup \textit{user friendly} semua sih ya Mbak, udah jelas semua gitu. Enggak ada yang apa ya, yang menyulitkan gitu ya Mbak. Gitu sih.

\textbf{Farah:} Kalau ada hal yang ingin diubah atau diperbaiki apa ya?

\textbf{Zufar:} Mungkin saya boleh liat lagi ya?

\textbf{Farah:} Boleh Mas, boleh, dieksplor lagi boleh Mas, saya tunggu juga boleh Mas.

\textbf{Zufar:} Paling menurut saya sebenarnya udah untuk peletakan lokasi udah bagus banget ya Mbak. Paling ini aja sih Mbak, \textit{font}-nya ya paling. Biasanya kalau seingat saya tuh biasanya lebih gede gitu loh, gitu biar lebih gampang dipencet gitu.

\textbf{Farah:} Oh oke Mas, makasih, sangat membantu untuk \textit{feedback}-nya. Oke, mungkin \textit{last question}. Tadi soal \textit{font}, mungkin ada lagi yang lain?

\textbf{Zufar:} Enggak ya, \textit{font} ini aja bagian navigasinya aja sih Mbak yang kurang gede gitu menurut saya.

\textbf{Farah:} Oke. Untuk yang fitur yang paling disuka apa Mas?

\textbf{Zufar:} Yang tadi AI sih Mbak. Soal AI sih menurut saya. Ini saya boleh lihat lagi ya sini. Soalnya tadi si bahasanya kan berdasarkan jawaban saya ya Mbak. Jadi soalnya kalau kita ngerjain, maksudnya kan sebenarnya latihan-latihan soal gini udah biasa ya Mbak, udah banyak gitu. Cuman kan pasti pembahasannya tuh ngikutin dari berdasarkan soalnya gitu. Sedangkan ini karena ada AI-nya jadi berdasarkan jawaban kita gitu sih pembahasannya.

\textbf{Farah:} Oke, mantap. Oke terima kasih Mas. Saya rasa untuk pengujiannya cukup sampai sini. Ini sangat berarti untuk menyelesaikan tugas akhir saya. Selamat malam Mas.

\textbf{Zufar:} Ya, selamat malam juga Mbak.

\section{Transkrip Wawancara Responden 2 (Amira Ghina Nurfansepta)}

\subsubsection*{1. Pendahuluan \& Persiapan}

\textbf{Farah:} Halo, terima kasih sudah bersedia meluangkan waktu Anda. Perkenalkan nama saya Farah Aulia. Biasa dipanggil Farah. Saya dari Institut Teknologi Bandung semester 8. Saat ini saya sedang melakukan pengujian untuk tugas akhir saya terkait aplikasi yang bernama Teman Belajar Algoritma Struktur Data. Nah, yang penting untuk digarisbawahi, yang sedang diuji adalah aplikasinya bukan kamu. Jadi, kalau nanti ada bagian yang membingungkan atau sulit ditemukan, itu justru masukan yang saya butuhkan, bukan kesalahan. Sesi ini akan berlangsung sekitar 15 sampai 30 menit, terdiri dari kamu mengerjakan beberapa tugas di aplikasi lalu mengisi kuesioner singkat. Kemudian selanjutnya saya akan mengamati dan mencatat tapi saya tidak akan membantu selama pengerjaan tugas. Kalau kamu benar-benar \textit{stuck}, boleh bilang "saya menyerah" dan kita lanjut ke tugas berikutnya. Sebelum kita mulai, apakah ada pertanyaan?

\textbf{Amira:} Tidak ada.

\textbf{Farah:} Oke bagus. Nah, selanjutnya ee silakan buka \textit{website} yang saya kirim di \textit{chat} dan silakan \textit{share screen}.

\textbf{Amira:} Oke. Eh... \textit{share screen}. Hem...

\textbf{Farah:} Bentar ya.

\textbf{Amira:} Aman. Silakan boleh e... mengambil waktu tidak masalah. Ya, ini karena ada setingan... Nah.

\textit{(Amira berhasil membagikan layarnya)}

\textbf{Amira:} Baik, apakah sudah terlihat layar saya?

\textbf{Farah:} Iya, sudah terlihat. Nah, sebelum masuk ke aplikasi, Anda bisa melakukan pendaftaran akun terlebih dahulu. Silakan lakukan pendaftaran.

\textbf{Amira:} Baik.

\textit{(Amira mulai mengisi form pendaftaran)}

\textbf{Amira:} Untuk email-nya, sesuai dengan email saya atau?

\textbf{Farah:} E terserah, boleh email bebas karena daftar dan masuk tidak termasuk dalam kriteria pengujian.

\textbf{Amira:} Baik.

\bigskip

\subsubsection*{2. Sesi Pengujian Aplikasi (Usability Testing)}

\textbf{Farah:} Oke, baik. Ini adalah halaman eee pertama yang Anda bisa akses dari Teman Belajar Algoritma dan Struktur Data. Nah, kali ini kita coba mulai Anda membuka eeee bagian topik dan pilih eeee \textit{list} linear untuk dipelajari.

\textbf{Amira:} Yang ini?

\textbf{Farah:} Eee \textit{list} linear.

\textbf{Amira:} \textit{List} linear. Oh \textit{sorry, sorry}, yang ini.

\textbf{Farah:} Iya. Nah, ee bagaimana pendapat Anda terkait navigasi ke halaman topik?

\textbf{Amira:} Sudah cukup baik, karena terdapat di sebelah atas ini, jadi kita bisa melihatnya dengan mudah, menu-menu apa saja yang ada. Dan se... ketika memilih topik, kita tinggal mem... melihat topik apa aja yang ada, kita \textit{scrolling}. Lalu kita pencet \textit{mulai} dan kita langsung masuk ke dalam materi dari topik-topik yang ada.

\textbf{Farah:} Oke, baik. Terima kasih atas ee tanggapannya. Nah, sekarang lakukan tugas untuk menyelesaikan materi. Jadi, baca seluruh materi dan tekan tombol eee selesaikan materi yang ada.

\textit{(Amira membaca materi dan menekan tombol)}

\textbf{Farah:} Oke, baik, terima kasih. Nah, berikutnya masuk ke halaman contoh dan baca contoh, serta selesaikan eee contoh tersebut.

\textit{(Amira membaca contoh dan menekan tombol)}

\textbf{Farah:} Oke. Sebelum kita lanjut ke latihan, apakah ada tanggapan terkait materi dan contoh yang ada?

\textbf{Amira:} Menurut saya sudah sangat baik, ya, karena ini, saya cukup suka karena berwarna, tidak hitam-putih saja di sini. Ada warna-warna yang bisa memanjakan mata saya.

\textbf{Farah:} Oke, baik. Terima kasih. Selanjutnya masuk ke halaman latihan. Nah, latihan ini adalah soal yang di-\textit{generate} oleh AI berdasarkan materi yang telah Anda pelajari. Jadi, membutuhkan waktu untuk men-\textit{generate}-nya. Silakan coba kerjakan satu soal saja dan eee baca hasil evaluasi yang diberikan oleh AI. Eee Anda boleh e... menjawab setahu Anda, tidak perlu benar.

\textbf{Amira:} Baik.

\textit{(Amira mengerjakan soal latihan)}

\textbf{Amira:} Langsung dinilai, soal?

\textbf{Farah:} Iya, boleh.

\textit{(Amira membaca hasil evaluasi)}

\textbf{Farah:} Oke, baik. Di sini Anda bisa melihat nilainya dan rincian penilaian yang telah dievaluasi oleh AI. Silakan baca evaluasinya dan berikan pendapat Anda terkait latihan dan evaluasi yang diberikan oleh AI ini.

\textbf{Amira:} Ya, di sini saya memang tidak menjawab dengan... lengkap, karena saya hanya menjelaskan perbedaan antara \textit{list} biasa dan \textit{double linked list}. Sementara soalnya itu, soal yang diberikan terdapat... saya disuruh untuk eee menyertakan kelebihan dan kekurangan dari masing-masing eee tipe data ini... struktur data tersebut. Jadi, ya memang jawaban saya belum lengkap, dan eee di sini terdapat eeee apa yang perlu saya tambahkan, apa yang kurang dari jawaban saya, dan ya, ini sudah benar, sih.

\textbf{Farah:} Oke, baik. Eee coba tandai selesai untuk bagian latihan dan silakan beralih ke ringkasan dan selesaikan bagian ringkasan.

\textit{(Amira menekan tombol dan beralih ke halaman ringkasan)}

\textbf{Farah:} Oke, baik. Jadi, eee bagaimana tanggapan Anda terkait bagian ringkasan?

\textbf{Amira:} Ya, sekali lagi saya suka apabila terdapat banyak warna, karena saya orangnya visual. Jadi, saya suka apabila warna-warna dan juga ini eee dikelompokkan berdasarkan topik-topik atau materi-materi yang eee sudah jelas seperti ini. Struktur berkait, representasi, variasi \textit{list} linear, \textit{stack, queue}, dan lain-lainnya.

\textbf{Farah:} Oke, terima kasih. Selanjutnya lakukan \textit{task} untuk melihat progres belajar Anda. Jadi, silakan Anda navigasi ke halaman progres dan lihat penyelesaian Anda terkait \textit{list} linear.

\textit{(Amira membuka halaman progres belajar)}

\textbf{Farah:} Oke. Nah, bisa Anda lihat bahwa progres telah selesai 100\% karena Anda telah menyelesaikan materi, latihan, contoh, dan ringkasan. Oke. Karena semua \textit{task} sudah berhasil Anda selesaikan dengan baik, selanjutnya kita masuk ke kuesioner. Jadi, sudah saya kirimkan kuesioner Google Form-nya ada di WhatsApp. Silakan isi tanpa saya mengetahui eee isi Anda. Jadi, silakan Anda \textit{stop share screen} sekarang.

\textbf{Amira:} Baik, saya isi sekarang ya.

\textbf{Farah:} Iya, terima kasih.

\bigskip

\subsubsection*{3. Sesi Kuesioner (SUS) \& Debriefing}

\textit{(Jeda saat Amira mengisi kuesioner)}

\textbf{Amira:} Baik, sudah saya \textit{submit} ya, untuk kuesionernya.

\textbf{Farah:} E baik, terima kasih banyak. Nah, berikutnya kita masuk ke sesi eee tambahan. Jadi eee saya ingin bertanya terkait pertanyaan-pertanyaan umum. Bagaimana kesan Anda terkait aplikasi ini?

\textbf{Amira:} Menurut saya, kesan saya, uh, menggunakan aplikasi ini... karena menurut saya Algoritma dan Struktur Data itu merupakan mata kuliah yang cukup sulit, ya, karena pada saat saya berkuliah dulu, eee cukup banyak teman-teman saya, dan juga saya sendiri merasa kesulitan dalam eee mengerjakan dan mempelajari, dan juga memahami materi-materi yang eee diberikan. Dengan adanya eee aplikasi ini atau \textit{website} ini, saya rasa se... eeee semakin mudah bagi saya untuk memahaminya.

\textbf{Farah:} Oke, baik. Eee terima kasih atas kesannya. Selanjutnya ee saya ingin bertanya, \textit{part} mana yang menurut Anda paling membingungkan?

\textbf{Amira:} Membingungkan... saya rasa nggak ada ya, karena semuanya sudah sangat jelas di situ. Jadi eee menurut saya se... sudah sangat mudah untuk dinavigasi, karena di situ sudah ada keterangan yang lengkap, jadi ya bagi saya tadi saya tidak mengalami kebingungan sama sekali, sih, dalam mengoperasikan eee aplikasinya.

\textbf{Farah:} Mmm, \textit{I see}. Oke, berikutnya eee apa bagian yang menurutmu paling membantu atau paling kamu suka?

\textbf{Amira:} Bagian yang paling membantu bagi saya yaitu ketika AI memberikan saya pertanyaan atas materi yang ada. Sehingga saya bisa me... lakukan, atau menguji diri saya, apakah saya sudah benar-benar memahami materinya, dan ketika saya memberikan jawaban eee di situ juga saya mendapatkan eeee evaluasi atau apakah jawaban saya ini sudah cukup benar atau belum. Jadi dari situ saya bisa me... apa ya... membuat saya semakin yakin lah, dalam belajar Algoritma dan Struktur Data ini.

\textbf{Farah:} Oke, baik. Terima kasih atas eee hal yang paling membantu. Selanjutnya ini pertanyaan terakhir, kalau ada satu hal yang di... kalau ada hal yang bisa diubah dari aplikasi ini, itu apa?

\textbf{Amira:} Mmm... apa ya... Menurut saya eee... tapi ini tidak eee... besar, ya, tapi mungkin dalam waktu apa... dalam menunggu soal untuk di-\textit{generate} AI, tapi memang saya sadari untuk \textit{generate} AI... apa, AI butuh waktu untuk men-\textit{generate} soal, tapi ya memang... ya itu. Jadi saya sebagai \textit{user}, menunggu, gitu.

\textbf{Farah:} Oh, saya mengerti. Oke, baik. Eee saya kira cukup untuk pengujian kali ini. Eee sudah... saya berterima kasih telah membantu saya dalam mengerjakan tugas akhir saya, semoga saya cepat lulus. Terima kasih.

\textbf{Amira:} Sama-sama.

\section{Transkrip Wawancara Responden 3 (Muhammad Faiz Fahri)}

\subsubsection*{Sesi Pembukaan (Pengantar \& Instruksi)}

\textbf{Farah:} Halo, perkenalkan nama saya Farah. Terima kasih sudah bersedia meluangkan waktu. Kita sedang menguji sebuah aplikasi bernama teman belajar algoritma dan struktur data. Nah, ini adalah bagian dari tugas akhir saya. Yang penting untuk digarisbawahi, yang sedang diuji itu adalah aplikasinya, bukan Anda. Jadi jika Anda mengalami kesulitan, tidak masalah, justru itu menjadi masukan untuk saya. Penelitian ini akan berjalan selama 15 sampai 30 menit. Selama mengerjakan, silakan kerjakan sesuai dengan instruksi saya. Saya akan mengamati dan mencatat, tapi saya tidak akan membantu. Jadi ketika Anda mengalami kesulitan, Anda bisa bilang saja. Apakah ada pertanyaan?

\textbf{Faiz:} Sudah jelas.

\subsubsection*{Sesi Pengujian Aplikasi (Task)}

\textbf{Farah:} Oke, baik, terima kasih. Selanjutnya kita bisa \textit{share screen} dan Anda bisa menggunakan aplikasi ini. Untuk tes pertama, silakan daftar akun terlebih dahulu.

\textbf{Faiz:} Baik. Daftar akun... \textit{[Faiz melakukan proses pendaftaran]}

\textbf{Farah:} Oke, baik. Jadi tadi untuk halaman \textit{login} dan daftar itu tidak termasuk ke dalam kriteria dalam penelitian ini. Jadi \textit{email} yang diverifikasi itu tidak benar sepenuhnya tidak masalah. Oke, tes berikutnya adalah, silakan masuk ke halaman topik dan pilih topik \textit{List Linear}.

\textbf{Faiz:} Oke. \textit{List Linear}. Iya.

\textbf{Farah:} \textit{List Linear}, oke. Silakan mulai belajar. Setelah Anda masuk ke halaman tersebut, silakan baca materinya dan tandai selesai. Tidak perlu dibaca semua, tidak masalah.

\textbf{Faiz:} \textit{[Membuka materi]} Oke, tandai selesai.

\textbf{Farah:} Iya. Apakah Anda mengalami kesulitan ketika melakukan ini?

\textbf{Faiz:} Tidak, ini di sebelah kiri sudah ada panduannya juga kan, mau... ini sedang di \textit{section} mana dan enggak ada banyak lagi.

\textbf{Farah:} Oke, baik, terima kasih. Selanjutnya silakan buka bagian contoh dan baca contoh, kemudian tandai selesai apabila sudah.

\textbf{Faiz:} Oke. \textit{[Membuka bagian contoh]} Oh, ini pembahasannya. Iya. Oke.

\textbf{Farah:} Oke, apakah Anda mengalami kesulitan atau ada masukan terkait fitur ini?

\textbf{Faiz:} Ini contoh, ya? Iya, contoh. Tadi bingung aja sih pas masuk contoh, ini enggak ngeh ini langsung soalnya. Ternyata ini langsung soal dan ini juga langsung ada pembahasannya juga.

\textbf{Farah:} Iya. Apakah ada masukan atau kendala, atau menurut Anda sudah oke?

\textbf{Faiz:} Kalau bisa ada penanda yang lebih jelas kalau misalkan ini udah langsung contoh. Ini kan kecil gini ya soalnya. Oh, \textit{I see}.

\textbf{Farah:} Oke, baik, terima kasih atas masukannya. Selanjutnya masuk ke bagian latihan dan coba kerjakan latihan yang ada. Soal ini berasal dari AI, jadi AI akan men-\textit{generate} soal ini. Anda bisa melakukan ganti soal.

\textbf{Faiz:} Ganti soal?

\textbf{Farah:} Iya. Dan AI akan men-\textit{generate} soal baru untuk Anda. Kemudian silakan kerjakan soal sesuai dengan apa yang Anda pahami, tidak perlu benar. Nanti baca evaluasi yang diberikan oleh AI.

\textbf{Faiz:} Oke. Kepanjangan ya, ganti soal aja. \textit{[Melakukan ganti soal dan mengerjakan]} Jelasin... Ini dinilai, kan?

\textbf{Farah:} Oke. Ini langsung dinilai dengan AI, ya.

\textbf{Faiz:} Oke. \textit{[Membaca hasil evaluasi]} Langsung sudah benar, oke. Tidak masalah. Jawaban sangat minim, oke.

\textbf{Farah:} Silakan dibaca dan tandai selesai apabila sudah, dan tolong berikan \textit{feedback} terkait evaluasi dan soal yang diberikan oleh AI, apakah ini membantu atau tidak.

\textbf{Faiz:} \textit{[Membaca evaluasi]} Tandai selesai.

\textbf{Farah:} Iya. Oke, jadi bagaimana pendapat Anda terkait fitur ini?

\textbf{Faiz:} Menurut saya ini sangat bagus untuk latihan soal dan mendalami \textit{knowledge} tentang apa yang sedang dipelajari, ya. Karena ini langsung dicelupin langsung dan kita disuruh kasih pemahaman kita sendiri tentang apa yang udah kita pelajari. Dan kalau misalkan salah, langsung dikoreksi, ya. Dan ini juga udah oke banget sih karena langsung di-\textit{point out} juga apa yang harus... bagian mana yang harus kita \textit{improve} dari jawaban kita gitu.

\textbf{Farah:} Oke, baik. Terima kasih atas \textit{feedback}-nya. Selanjutnya bagian ringkasan. Silakan baca bagian ringkasan dan tandai selesai apabila sudah.

\textbf{Faiz:} \textit{[Membuka bagian ringkasan]} Oke, baik.

\textbf{Farah:} Apa pendapat Anda terkait fitur ini?

\textbf{Faiz:} Menurut saya ringkasannya ini udah oke. Cuma mungkin ini kan langsung mendalam banget, ya. Mungkin bisa lebih apa... dari mendasar dulu. Dan di sini saya juga tidak melihat ada visualisasi. Menurut saya untuk materi ini akan lebih mudah dipahami jika ada visualisasi. Tapi di sini yang saya lihat hanya teks dan banyak istilah-istilah asing yang mungkin orang awam tidak ketahui gitu.

\textbf{Farah:} Oh, mengerti. Iya. Terima kasih atas masukannya. Selanjutnya beralih ke halaman progres dan lihat progres belajar Anda terkait \textit{List Linear}.

\textbf{Faiz:} 10\%. Progres. Oke.

\textbf{Farah:} Oke, bisa Anda lihat bahwa progres Anda telah 100\%, berarti Anda telah menyelesaikan materi \textit{List Linear}. Untuk pengujian seluruh \textit{task}-nya sudah bisa berhasil Anda selesaikan. Selanjutnya adalah pengisian formulir. Jadi saya akan kirimkan formulir kepada Anda dan Anda bisa isi terlebih dahulu.

\textbf{Faiz:} \textit{[Mengisi formulir evaluasi]} Oke, baik.

\textbf{Farah:} Terima kasih untuk pengisian form-nya.

\subsubsection*{Sesi Wawancara (Post-Test) \& Penutup}

\textbf{Farah:} Selanjutnya, ini adalah pertanyaan tambahan. Jadi saya ingin bertanya lebih lanjut. Secara umum, bagaimana kesan Anda menggunakan aplikasi ini?

\textbf{Faiz:} Kesan saya, ini aplikasi menarik, ya. Ini aplikasi teman belajar untuk belajar algoritma dan struktur data, khususnya untuk mahasiswa-mahasiswa informatika yang dapet mata kuliah ini. Menurut saya ini aplikasi yang bakalan bisa membantu mahasiswa-mahasiswa tersebut. Khususnya dari bagian \textit{engagement}, ya. Itu karena langsung ada bagian pertanyaannya yang bisa langsung menguji mahasiswa tersebut dan bisa langsung membantu mahasiswa tersebut kurangnya di bagian mana gitu.

\textbf{Farah:} Oke, baik, terima kasih. Selanjutnya, untuk hal-hal yang perlu diperbaiki dari aplikasi ini, apa saja, ya?

\textbf{Faiz:} Pertama menurut saya di bagian visualisasi, ya. Saya rasa untuk materi struktur data khususnya itu akan lebih mudah dipahami jika ada banyak visualisasi, khususnya tadi di \textit{linked list}, ya. Tadi saya kurang melihat visualisasinya.

\textbf{Farah:} Oke. Apa bagian yang menurut Anda paling menarik atau yang paling Anda suka?

\textbf{Faiz:} Menurut saya yang paling menarik di fitur AI-nya, ya. Tadi misalkan kita kurang puas dengan pertanyaannya, atau jika merasa pertanyaannya kegampangan nih, kita bisa langsung ganti soal ke soal yang menurut kita \textit{challenging} gitu. Di bagian yang tidak sukanya, menurut saya di bagian penjelasan, sih, ya. Di bagian penjelasan itu menurut saya udah langsung terlalu mendalam secara bahasa, dan sebenarnya secara visual itu tidak ada masalah. Cuman masalah di kontennya aja, itu belum dari dasar, ya. Berarti harusnya benar.

\textbf{Farah:} Oke, baik. Pertanyaan terakhir, kalau ada yang bisa diubah dari aplikasi ini, apa yang bisa diubah?

\textbf{Faiz:} Menurut saya pertama, adakan fitur \textit{night mode}.

\textbf{Farah:} Oke.

\textbf{Faiz:} Dan yang kedua, saya akan tambahkan visualisasi aja sih.

\textbf{Farah:} Oke, baik. Terima kasih banyak sudah membantu saya dalam mengerjakan tugas akhir ini. Semoga saya bisa lulus tepat waktu.

\textbf{Faiz:} Amin. Amin.

\textbf{Farah:} Terima kasih banyak, Faiz. Selamat malam.

\textbf{Faiz:} Oke, selamat malam.

\section{Transkrip Wawancara Responden 4 (Thea)}

\subsubsection*{1. Pendahuluan \& Persiapan}

\textbf{Farah:} Halo, selamat malam! Perkenalkan nama saya Farah. Terima kasih sudah bersedia meluangkan waktu Anda. Saya sedang menguji sebuah aplikasi bernama teman belajar algoritma dan struktur data. Ini adalah bagian dari tugas akhir saya. Yang penting untuk digarisbawahi, yang sedang diuji adalah aplikasinya, bukan kamu. Jadi kalau nanti ada bagian yang membingungkan atau sulit ditemukan, itu justru menjadi masukan bagi saya, dan bukan kesalahan Anda. Nah, saat mengerjakan, saya akan mengamati dan mencatat, tapi saya tidak akan membantu selama pengerjaan. Itu menjadi bagian dari pengujian. Kalau kamu benar-benar \textit{stuck}, kamu boleh bilang untuk menyerah dan kita akan lanjut ke tugas berikutnya. Apakah ada pertanyaan sebelum kita mulai?

\textbf{Thea:} Aman-aman. Nggak ada.

\textbf{Farah:} Oke. Baik, terima kasih. E silakan buka \textit{website} yang saya kirim di WhatsApp sebelumnya dan silakan lakukan \textit{share screen} untuk \textit{website} tersebut.

\textbf{Thea:} Udah buka.

\textbf{Farah:} E silakan \textit{share screen}.

\textbf{Thea:} Oh iya ya. Udah. Udah keliatan nggak?

\textbf{Farah:} Oke, sudah terlihat. E silakan lakukan pendaftaran terlebih dahulu.

\bigskip

\subsubsection*{2. Sesi Pengujian Aplikasi (Usability Testing)}

\textit{(Thea mencoba mendaftar/login)}

\textbf{Thea:} Loh, ini kenapa Far?

\textbf{Farah:} Emm... Coba kata sandinya.

\textbf{Thea:} Oke, oke, oke.

\textbf{Farah:} Oke, baik tidak masalah. Jadi sebenarnya daftar dan \textit{login} itu tidak tidak termasuk dalam pengujian saya, jadi e itu tidak termasuk ke dalam e pengujian untuk itu. Nah, mari kita mulai untuk pengujiannya. Tugas pertama yang perlu Anda lakukan adalah buka halaman topik dan pilih \textit{list} linear, kemudian pelajari topik tersebut.

\textit{(Thea membuka topik dan membaca)}

\textbf{Farah:} Oke, tandai selesai. He eh. Apakah ada masukan terkait materi e \textit{list} ini?

\textbf{Thea:} E... bentar aku liat sambil liat-liat lagi.

\textbf{Farah:} Boleh.

\textbf{Thea:} Dari aku aman aja sih. Bentar bentar, aku mau tanya, aku mau tanya. Ringkasan empat alternatif tuh dia maksudnya emang ringkasan gitu atau gimana?

\textbf{Farah:} Oke, baik. Untuk ringkasan empat alternatif itu adalah e alternatif dari implementasi \textit{list}, jadi bisa secara implisit, eksplisit, atau e dan lainnya. Itu bentuk \textit{list} aja sih. Jadi bukan ringkasan dari e ringkasan semuanya.

\textbf{Thea:} Oh, oalah. Oke oke oke. Aman aja sih kalau aku.

\textbf{Farah:} Oke baik. Selanjutnya e silakan buka halaman contoh dan pelajari contoh, kemudian tandai selesai apabila sudah.

\textit{(Thea membuka halaman contoh)}

\textbf{Thea:} Udah.

\textbf{Farah:} Apakah ada masukan atau \textit{feedback} terkait dengan bagian contoh?

\textbf{Thea:} Mmm... Sebenarnya nggak ada sih, aman-aman aja sih kalau aku.

\textbf{Farah:} Oke baik, terima kasih. Berikutnya lanjut ke halaman latihan, kerjakan satu soal latihan. Jadi ini menunggu karena soalnya ada e di-\textit{generate} oleh AI. Kerjakan satu soal dan baca evaluasinya, kemudian tandai selesai apabila sudah. Anda tidak perlu e jawaban yang benar, masukkan jawaban yang e Anda tahu saja tidak masalah.

\textbf{Thea:} "Jelaskan definisi ADT." Tidak perlu benar, tidak masalah. Oh, oke oke oke. Ya ya, saya mau benar juga nggak apa-apa.

\textbf{Farah:} \textit{(Tertawa)}

\textbf{Thea:} Nyontek sini aja, oh nggak bisa ya. Bisa.

\textit{(Thea mengerjakan soal latihan)}

\textbf{Thea:} Udah. Ini \textit{next} aja.

\textbf{Farah:} Ingat tolong berikan \textit{feedback} dari soal yang telah Anda kerjakan, apakah Anda setuju dengan evaluasi yang diberikan.

\textbf{Thea:} Ini kenapa belum dijawab?

\textbf{Farah:} Oh, itu karena AI mengevaluasi bahwa jawabannya kurang tepat.

\textbf{Thea:} Oh. Kayaknya kalau aku tampilannya aja sih, kan ini jawaban aku ya, kotakannya kan abu, tapi yang bawahnya tuh jawabannya dia, "jawaban tidak lengkap", itu kotakannya juga abu. Kayak... mungkin dibedain gitu.

\textbf{Farah:} Oh, benar. Terima kasih atas masukannya, Thea. Oke, baik. Apakah ada masukan yang lainnya?

\textbf{Thea:} E... Oh ini tuh di-\textit{generate} AI-nya atau gimana ya?

\textbf{Farah:} Iya.

\textbf{Thea:} Oh kalau bisa dia suruh ngambilnya yang dari ini sih, nama-nama \textit{tab}-nya ini.

\textbf{Farah:} Oke, terima kasih, menarik. Oke, silakan tandai selesai apabila tidak ada \textit{feedback} lagi. Selanjutnya silakan baca e bagian ringkasan dan selesaikan bagian ringkasan apabila dirasa sudah cukup.

\textit{(Thea membuka bagian ringkasan)}

\textbf{Thea:} Udah.

\textbf{Farah:} Oke. E apakah ada masukan terkait bagian ringkasan?

\textbf{Thea:} E... nggak ada sih. Cuman aku bingungnya tuh kayak... kenapa ringkasan tuh baru bisa diakses setelah latihan gitu. Eh, iya nggak sih?

\textbf{Farah:} Iya, iya.

\textbf{Thea:} Aku mikirnya kenapa nggak setelah baca materi gitu juga ada ringkasannya gitu.

\textbf{Farah:} Oke. Terima kasih atas masukannya. Jadi sebenarnya e ini sengaja dibuat untuk e alurnya seperti itu, jadi dari materi, contoh, latihan, kemudian ringkasan. Jadi e peserta bisa me-\textit{review} semua yang telah dia pelajari dan dia kerjakan di ringkasan. Tapi mungkin menjadi masukan untuk ringkasannya dipindah setelah e bagian materi atau contoh.

\textbf{Thea:} Oke oke.

\textbf{Farah:} Oke. Selanjutnya \textit{task} berikutnya silakan buka halaman progres dan lihat progres belajar Anda untuk bagian \textit{list}.

\textit{(Thea melihat halaman progres)}

\textbf{Thea:} Oke.

\textbf{Farah:} Oke. E sudah. Saya rasa untuk tugasnya sudah selesai sampai sini. Selanjutnya kita masuk ke halaman e pengisian \textit{form}. Jadi Anda bisa \textit{stop share screen} dan isi \textit{form} yang telah saya kirimkan di WhatsApp.

\bigskip

\subsubsection*{3. Sesi Kuesioner (SUS) \& Debriefing}

\textbf{Thea:} Oke. Yang SUS ini kan ya?

\textbf{Farah:} Iya. Itu hanya ada 10 soal.

\textit{(Jeda saat Thea mengisi kuesioner SUS)}

\textbf{Thea:} Udah.

\textbf{Farah:} Oke baik e terima kasih. Ini ada tiga pertanyaan tambahan saja e untuk menutup sesi ini. Nah, e pertanyaan pertama, bagaimana kesan keseluruhan Anda terkait aplikasi ini?

\textbf{Thea:} E kesan keseluruhan ngebantu sih. Soalnya ini kan dari sini kita bisa ngelihat berbagai jenis topik yang ada, terus ada materinya itu udah terstruktur gitu secara langsung. Terus sama kayak mungkin yang paling bantu tuh apa, latihan itu ya. Soalnya kan dia pakai AI. Nah itu latihannya tuh bisa macam-macam gitu, nggak bakal bisa beda gitu. Jadi bakal bantu untuk belajar.

\textbf{Farah:} Baik, terima kasih. Kemudian untuk e apa fitur atau e \textit{task} yang paling menarik untuk Anda?

\textbf{Thea:} E latihan.

\textbf{Farah:} Oke. Apabila ada yang bisa diubah dari aplikasi ini, apa yang bisa diubah?

\textbf{Thea:} E kayaknya tambahin gambar aja sih biar kayak lebih menarik gitu, atau mungkin bisa ada kayak animasi kayak untuk menunjukkan proses insersinya dalam \textit{list} misalnya gitu, atau kayak proses \textit{delete}. Tapi kalau pakai ini juga udah jelas sih sebenarnya.

\textbf{Farah:} Oke baik. E saya rasa cukup sekian e pengujian dari saya. Terima kasih sudah membantu dalam pengujian tugas akhir saya. Semoga saya bisa lulus tepat waktu. Selamat malam, Thea.

\textbf{Thea:} Amin.

\section{Transkrip Wawancara Responden 5 (Vieri)}

\textit{Catatan: beberapa baris percakapan informal yang tidak relevan dengan hasil pengujian telah dihilangkan dari transkrip berikut, ditandai dengan keterangan dalam tanda kurung siku.}

\subsubsection*{1. Pendahuluan \& Persiapan}

\textbf{Farah:} Selamat malam, terima kasih sudah bersedia meluangkan waktu. Perkenalkan, nama saya Farah Aulia. Saya sedang menguji sebuah aplikasi bernama Teman Belajar Algoritma dan Struktur Data, sebagai bagian dari tugas akhir saya. Nah, yang penting untuk digarisbawahi, yang sedang diuji adalah aplikasinya bukan kamu. Jadi kalau nanti ada bagian yang membingungkan atau sulit ditemukan, itu menjadi masukan yang saya butuhkan, bukan kesalahan kamu. Sesi ini akan berlangsung selama 10 sampai 20 menit tergantung situasinya. Kamu nanti akan diperintahkan untuk mengerjakan beberapa tugas di aplikasi ini lalu mengisi kuesioner singkat. Nah, selama mengerjakan tugas, saya akan mengamati dan mencatat, tapi saya tidak akan membantu selama pengerjaan tugas. Itu memang bagian dari pengujian. Kalau kamu benar-benar \textit{stuck}, boleh bilang "saya menyerah" dan kita lanjut ke tugas berikutnya. Ada pertanyaan sebelum kita mulai?

\textbf{Vieri:} Cukup.

\textbf{Farah:} Oke, mantap. Nah, selanjutnya silakan \textit{share screen} ke halaman \textit{website} yang sudah saya kirimkan di \textit{chat}.

\textbf{Vieri:} Oke terus.

\textbf{Farah:} Terima kasih, sudah terlihat. Sekarang silakan lakukan pendaftaran di halaman \textit{website} tersebut.

\textit{(Vieri melakukan pendaftaran)}

\textbf{Farah:} Oke baik. Halaman \textit{login} dan daftar tidak termasuk ke dalam pengujian dari \textit{task} ini. Selanjutnya kita akan fokus ke \textit{task} yang pertama, yaitu navigasi ke halaman topik dan pilih topik \textit{list} linear untuk dipelajari.

\bigskip

\subsubsection*{2. Sesi Pengujian Aplikasi (Usability Testing)}

\textbf{Vieri:} Kenapa ya harus \textit{list} linear?

\textbf{Farah:} \textit{List} linear adalah salah satu contoh topik yang sengaja saya pilih untuk demonstrasi pada pengujian ini.

\textbf{Vieri:} Mmm. Enggak boleh milih ya?

\textbf{Farah:} Ya, silakan pilih \textit{list} linear dan mulai untuk belajar. Oke, baik. Selanjutnya masuk ke materi dan baca materinya, kemudian tandai selesai setelah materi selesai dibaca.

\textbf{Vieri:} Ini aku boleh nanya enggak?

\textbf{Farah:} Apa? Silakan, silakan boleh bertanya.

\textbf{Vieri:} Saya lagi males baca soalnya.

\textbf{Farah:} Materi tidak termasuk ke dalam penilaian, jadi silakan baca sekilas tidak masalah.

\textbf{Vieri:} Aduh. Susahnya ternyata belajar ASD tuh. \textit{Double linked list}. Mmm, terus nama \textit{next}. Hm. Ini apaan ya \textit{stack} ya. Bisa tolong jelasin enggak kenapa ada ini, \textit{stack} dengan bentuk struktur berkait?

\textbf{Farah:} Baik, terima kasih atas pertanyaannya. Jadi ada bertanya tentang \textit{stack} dengan \textit{list} linear ya? Nah, itu nanti jadi \textit{stack} yang diimplementasikan dengan \textit{list} linear. Jadi kita bisa tahu elemen teratasnya dengan bentuk \textit{list} linear. Apa yang Anda kurang paham?

\textbf{Vieri:} Oh, jadi ini sebenarnya \textit{stack}. Cuman \textit{value} dari \textit{stack}-nya tuh \textit{list} linear gitu ya?

\textbf{Farah:} Sepertinya seperti itu.

\textbf{Vieri:} Ada \textit{queue} juga. Oh. Aneh nih ada \textit{priority queue} pakai \textit{list} linear, susah dong. Oke ya udahlah tandai selesai. Terus.

\textbf{Farah:} Oke baik. Bagaimana pendapat Anda terhadap materi ini?

\textbf{Vieri:} Ya kalau pendapat saya sih... sebenarnya kalau tentang apa sih \textit{list} linear itu enggak perlu sampai jelasin yang \textit{queue} terus \textit{stack} sama \textit{priority queue} itu ya. Itu mungkin bisa dijelasin di materi lain yang lebih spesifik. Kayak contohnya yang materi-materi \textit{stack}, \textit{queue}, sama \textit{priority queue} itu bisa kamu jelasin ke materi yang LIFO sama FIFO.

\textbf{Farah:} Baik, terima kasih atas tanggapannya. Lanjut ke fitur berikutnya, silakan akses halaman contoh dan baca contohnya. Kemudian tandai selesai ketika sudah.

\textbf{Vieri:} Mmm. Oke. Udah. Halo bro.

\textbf{Farah:} Oke nah. Tandai selesai ketika sudah. Nah, bagaimana pendapat Anda terkait fitur contoh ini?

\textbf{Vieri:} Oke contohnya cukup jelas. He eh. Oke cukup jelas. Enggak ada masalah contohnya.

\textbf{Farah:} Baik selanjutnya, akses halaman latihan dan kerjakan salah satu soal yang ada di situ. Dan baca evaluasinya, tandai selesai ketika sudah.

\textbf{Vieri:} Kok lama ya? Ada \textit{bug} kah? Kayaknya boleh pakai AI enggak?

\textbf{Farah:} Silakan dikerjakan sesuai pemahaman Anda. Tidak perlu benar.

\textbf{Vieri:} Bentar ya saya mikir dulu. Ini kalau saya enggak tahu gimana?

\textbf{Farah:} Enggak apa-apa, tidak perlu benar.

\textbf{Vieri:} Enggak apa-apa ya. Udah. Terus. Halo.

\textbf{Farah:} Oke, silakan nilai soal ini dan baca hasil evaluasinya.

\textbf{Vieri:} Boleh saya lihat untuk nilainya? 60.

\textbf{Farah:} Oke silakan baca hasil evaluasi yang diberikan, apakah sesuai dan apakah mudah dipahami dan silakan berikan masukan apabila kurang.

\textbf{Vieri:} Mmm mungkin ini sih. Apa? Enggak enggak menyinggung kompleksitas tapi kenapa ada kompleksitasnya.

\textbf{Farah:} Baik, terima kasih jadi Anda... menjawab bahwa hasil evaluasi yang diberikan kurang tepat atau bagaimana?

\textbf{Vieri:} Ya kurang tepat untuk penilaiannya.

\textbf{Farah:} Oke, baik. Apakah ada hal yang lain?

\textbf{Vieri:} Udah cukup.

\textbf{Farah:} Oke selanjutnya silakan tandai selesai dan masuk ke halaman ringkasan dan baca ringkasannya. Kemudian tandai selesai apabila sudah.

\textbf{Vieri:} He eh. Udah.

\textbf{Farah:} Oke, apakah ada masukan atau \textit{feedback} terkait fitur ini?

\textbf{Vieri:} Cukup.

\textbf{Farah:} Oke, setelah itu tugas yang terakhir, silakan masuk ke halaman progres dan lihat progres belajar Anda terkait \textit{list} linear.

\textbf{Vieri:} Hm. Mana? Lihat progres belajar Anda terkait \textit{list} linear. Oke.

\textbf{Farah:} Baik Anda telah menyelesaikan \textit{list} linear 100\% materi. Untuk pengujian dari \textit{task} yang perlu Anda kerjakan telah selesai.

\bigskip

\subsubsection*{3. Sesi Kuesioner (SUS) \& Debriefing}

\textbf{Farah:} Selanjutnya kita masuk ke pengisian kuesioner. Silakan akses \textit{link} kuesioner yang saya kirim di \textit{chat}, Anda tidak perlu \textit{share screen} untuk mengisinya. Setelah selesai, tolong beritahu saya.

\textbf{Vieri:} Sensitif banget nanya-nanya. Bentar ya saya mikir dulu. Ini kalau saya enggak tahu gimana?

\textbf{Farah:} Enggak apa-apa, tidak perlu benar.

\textit{[Beberapa baris percakapan informal yang tidak relevan dengan hasil pengujian dihilangkan.]}

\textbf{Farah:} Silakan isi sesuai dengan apa yang Anda yakini.

\textbf{Farah:} Oke, baik. Selanjutnya silakan tutup \textit{share screen}-nya dan kita akan mulai di sesi terakhir yaitu sesi tambahan. Bagaimana kesan Anda keseluruhan terhadap aplikasi ini?

\textbf{Vieri:} Oke sih.

\textbf{Farah:} Halo. Bagaimana kesan Anda terhadap aplikasi ini secara keseluruhan?

\textbf{Vieri:} Oke cuman soalnya ini aja, validasinya agak jelek aja sih menurutku. Sama kadang materinya yang tadi tuh enggak ini, enggak terlalu relevan tapi dimasukin. Sama ini sih mungkin bisa tambahin perbanyak contohnya aja tiap-tiap sub-materi dikasih contoh masing-masing, kayak \textit{single linked list} itu dikasih contohnya, \textit{double linked list} dikasih contohnya. Tadi contohnya satu doang.

\textbf{Farah:} Baik, terima kasih. Jadi menurut Anda apa hal yang paling Anda sukai dari aplikasi ini?

\textbf{Vieri:} \textit{Generate} soalnya.

\textbf{Farah:} Apa yang menurut Anda paling membingungkan?

\textbf{Vieri:} Validasi soalnya.

\textbf{Farah:} Kalau ada yang bisa diubah dari aplikasi ini, apa yang bisa diubah?

\textbf{Vieri:} Mungkin contohnya sama validasinya lebih oke lah. Enggak terlalu \textit{strict}. Validasinya bisa diganti manusia aja.

\textbf{Farah:} Oh baik, terima kasih. Saya rasa cukup sekian pengujian kali ini. Anda telah membantu saya mengerjakan tugas akhir saya, semoga tugas akhir saya selesai dengan baik dan bisa lulus tepat waktu.

\textbf{Vieri:} Amin.

\textbf{Farah:} Oke baik saya akhiri selamat malam. Hm.
